% Compile: pdflatex guy_notes.tex (twice for the contents).
% Source: notes from guy 18_11_26.pdf, all 21 pages.
\documentclass[11pt,a4paper]{article}
\usepackage[T1]{fontenc}
\usepackage[utf8]{inputenc}
\usepackage{lmodern}
\usepackage[margin=25mm]{geometry}
\usepackage{amsmath,amssymb,amsthm}
\usepackage{xcolor}
\usepackage[hidelinks]{hyperref}
\usepackage{microtype}
\allowdisplaybreaks
\setlength{\parindent}{0pt}
\setlength{\parskip}{5pt}
\newcommand{\eps}{\varepsilon}
\newcommand{\OO}{\mathcal O}
\DeclareMathOperator{\ord}{ord}
\DeclareMathOperator{\Null}{Null}
\DeclareMathOperator{\rank}{rank}
\DeclareMathOperator{\spanof}{span}
\newcommand{\src}[1]{\par\smallskip{\small\color{gray}Source PDF: page(s) #1.}\par}
\newcommand{\edit}[2]{\par\begingroup\small\color{blue!55!black}\textbf{#1.} #2\par\endgroup}
\newtheorem{definition}{Definition}[section]
\newtheorem{proposition}[definition]{Proposition}
\newtheorem{example}[definition]{Example}
\numberwithin{equation}{section}
\title{Asymptotic Methods\\\large Guy's Handwritten Notes: Edited LaTeX Transcription}
\author{Source notes by Guy}
\date{Source file dated 18 November 2026}
\begin{document}
\maketitle
\begin{abstract}
An edited mathematical transcription of all 21 pages of
\texttt{notes from guy 18\_11\_26.pdf}. The order of topics and worked
examples follows the source. English and layout have been cleaned up;
repeated algebra has occasionally been combined. Blue editorial notes
identify substantive corrections, deductions, added hypotheses, and
unresolved notation. Source-page markers support comparison with the original.
This is an edited study edition, not a diplomatic, symbol-for-symbol copy.
\end{abstract}
\tableofcontents
\section*{Reading this edition}
\addcontentsline{toc}{section}{Reading this edition}
The source's $\ord(g)$ notation means two-sided comparability in absolute
value (often written $\Theta(g)$); it does not mean asymptotic equivalence.
The symbol $\ll$ means little-$o$ when used in asymptotic comparisons.
Unless stated otherwise, $\eps\to0^+$ and large real parameters tend to
$+\infty$. Corrections are displayed in the main mathematics and explained
locally. Additional reasoning is labelled \emph{Reconstruction} or
\emph{Clarification}. An editorial register appears at the end.

\section{Euler's constant}
\src{1--2}
Define
\begin{equation}
 H_n=\sum_{j=1}^n\frac1j,\qquad a_n=H_n-\ln n.
\end{equation}
We claim that $a_n$ converges to Euler's constant $\gamma$.
\subsection{Existence of the limit}
For $j<t<j+1$,
\[
 \frac1{j+1}<\frac1t<\frac1j.
\]
Integrating on $[j,j+1]$ and summing gives, for $n>1$,
\[
 H_n-1<\ln n<H_{n-1},\qquad \frac1n<a_n<1.
\]
Furthermore,
\begin{align}
 a_{n+1}-a_n
 &=\frac1{n+1}-\ln(n+1)+\ln n\\
 &=\frac1{n+1}+\ln\left(1-\frac1{n+1}\right)<0.
\end{align}
Indeed $h(x)=\ln(1-x)+x$ satisfies $h(0)=0$ and
$h'(x)=-x/(1-x)<0$ for $0<x<1$.
Thus $(a_n)$ is decreasing and bounded below, so
\[
 \gamma=\lim_{n\to\infty}(H_n-\ln n)=0.57721\ldots.
\]
\edit{Correction (p.~1)}{The logarithmic inequality uses
$\ln(1-x)+x$, as confirmed by the derivative in the notes. The handwritten
claim appears to have a minus sign. Strict bounds above require $n>1$;
$a_1=1$.}

\subsection{Integral representations}
The geometric-sum identity yields
\[
 \sum_{j=1}^n q^j=\frac{q(1-q^n)}{1-q}.
\]
Since $1/j=\int_0^\infty e^{-jt}\,dt$,
\begin{equation}
 H_n=\int_0^\infty
 \frac{e^{-t}(1-e^{-nt})}{1-e^{-t}}\,dt.
\end{equation}
For fixed $n$, this integrand tends to $n$ as $t\to0^+$.
If $n\to\infty$ is taken first, the resulting integrand is
$e^{-t}/(1-e^{-t})\sim1/t$, which is not integrable at zero.
Similarly, interchanging nonnegative integrals,
\begin{align}
 \ln n&=\int_1^n\frac{dx}{x}
 =\int_1^n\int_0^\infty e^{-xt}\,dt\,dx\\
 &=\int_0^\infty\frac{e^{-t}-e^{-nt}}t\,dt.
\end{align}
Subtracting the representations before passing to the limit gives
\begin{equation}\label{eq:gammaint}
 \gamma=\int_0^\infty e^{-t}
 \left(\frac1{1-e^{-t}}-\frac1t\right)\,dt.
\end{equation}
The cancellation near zero follows from
\[
 1-e^{-t}=t-\frac{t^2}2+\OO(t^3),\qquad
 \frac1{1-e^{-t}}=\frac1t+\frac12+\OO(t).
\]
\edit{Correction and reconstruction (p.~2)}{The first-order Taylor signs
are corrected: $1-e^{-nt}=nt+\OO(t^2)$ for fixed $n$, and
$e^{-t}=1-t+\OO(t^2)$. The omitted justification of
\eqref{eq:gammaint} can be supplied as follows. The difference between
$a_n$ and the displayed limiting integral is
$\int_0^\infty e^{-nt}r(t)\,dt$, where
$r(t)=1/t-1/(e^t-1)$ is positive, continuous at zero with value $1/2$,
and bounded. Hence this difference is $\OO(1/n)$.}

\section{The Gamma function}
\src{2--4}
For $x>0$, define
\begin{equation}
 \Gamma(x)=\int_0^\infty e^{-t}t^{x-1}\,dt.
\end{equation}
In particular $\Gamma(1)=1$. Integration by parts gives
\[
 \Gamma(x+1)=[-e^{-t}t^x]_0^\infty
 +x\int_0^\infty e^{-t}t^{x-1}\,dt=x\Gamma(x).
\]
Consequently $\Gamma(n+1)=n!$ for integers $n\ge0$.
With $t=u^2$,
\[
 \Gamma\left(\tfrac12\right)
 =2\int_0^\infty e^{-u^2}\,du=\sqrt\pi.
\]
\edit{Clarification (p.~3)}{The informal expression $(-1/2)!$ means
$\Gamma(1/2)$; the substitution requires the factor $2$ above.}
For complex $z$ with $\Re z>0$,
\[
 \Gamma(z)=\int_0^\infty e^{-t}t^{z-1}\,dt,\qquad
 \Gamma(z+1)=z\Gamma(z).
\]
The recurrence continues $\Gamma$ meromorphically to the complex plane,
with poles at $0,-1,-2,\ldots$.
\edit{Clarification (p.~3)}{Continuation by $\Gamma(z)=\Gamma(z+1)/z$
does not assign finite values at the poles.}

\subsection{A binomial series}
For $p>0$, let $f(x)=(1+x)^{-p}$. Then
\[
 f^{(n)}(0)=(-1)^n p(p+1)\cdots(p+n-1)
 =(-1)^n\frac{\Gamma(p+n)}{\Gamma(p)}.
\]
Thus
\begin{equation}
 (1+x)^{-p}=\sum_{n=0}^\infty
 \frac{(-1)^n\Gamma(p+n)}{n!\Gamma(p)}x^n,\qquad |x|<1.
\end{equation}
\edit{Clarification (pp.~3--4)}{The convergence domain $|x|<1$ is added;
the zeroth derivative uses an empty product equal to $1$.}

\section{Order notation and asymptotic expansions}
\subsection{Order notation}
\src{4--5}
As $x\to0$, define:
\begin{align*}
 f=\OO(g)&\quad\Longleftrightarrow\quad
 \exists C,\delta>0:\ |f(x)|\le C|g(x)|\quad(0<|x|<\delta),\\
 f=o(g)&\quad\Longleftrightarrow\quad
 \forall C>0\ \exists\delta>0:\ |f(x)|\le C|g(x)|
 \quad(0<|x|<\delta),\\
 f=\ord(g)&\quad\Longleftrightarrow\quad
 \exists c,C,\delta>0:\ c|g(x)|\le|f(x)|\le C|g(x)|
 \quad(0<|x|<\delta).
\end{align*}
The function $g$ is called a \emph{gauge function}. If $g$ is nonzero
in a punctured neighbourhood, the definitions are equivalent to:
$f/g$ bounded, $f/g\to0$, and $|f/g|$ bounded above and away from zero,
respectively. Definitions at infinity are analogous.

Examples include
\begin{align*}
 \sin x&=x+\OO(x^3)=x+o(x)=x+\ord(x^3) &&(x\to0),\\
 x^2\ln x+x^3&=\begin{cases}\OO(x^3),&x\to\infty,\\
 \OO(x^2\ln x),&x\to0^+,\end{cases}\\
 x^2-1&=\ord(x-1)&&(x\to0),\\
 x^{-2}&=o(x^{-1})&&(x\to\infty),\\
 \sinh x&=\begin{cases}\OO(e^x),&x\to\infty,\\
 \OO(e^{-x}),&x\to-\infty.\end{cases}
\end{align*}
\edit{Clarification (p.~5)}{The somewhat unusual $x^2-1$ example is
retained with the source's limit $x\to0$; it is valid there, and also
as $x\to1$. The logarithmic limit is one-sided for real logarithms.}
For fixed $k>0$ and $\eta>0$,
\[
 (\ln x)^k\ll x^\eta,\qquad x^k\ll(1+\eta)^x,
 \qquad k^x\ll x^{\eta x}\qquad(x\to\infty).
\]
Although one writes $f=\OO(g)$, the notation describes membership
in a class of functions. It does not behave as ordinary equality.
For example, with
\[
 f(n)=\begin{cases}1,&n\text{ odd},\\n^2,&n\text{ even},\end{cases}
 \qquad g(n)=n^2,
\]
we have $f=\OO(g)$, but neither $f=o(g)$ nor $f=\ord(g)$.

\subsection{Arithmetic and common pitfalls}
\src{5--6}
For suitable functions near the stated limit,
\begin{align*}
 \OO(g)+o(g)&=\OO(g),\\
 \OO(g_1)+\OO(g_2)&\subseteq\OO(|g_1|+|g_2|),\\
 \OO(f)\,o(g)&\subseteq o(fg).
\end{align*}
Here class identities are understood as usual asymptotic shorthand.
For example, as $x\to0$,
\[
 [x+\OO(x^3)]^2=x^2+\OO(x^4)+\OO(x^6)=x^2+\OO(x^4),
 \qquad (2+\sin x)o(x^2)=o(x^2).
\]
Also $\OO(x^2)+o(x)+\OO(x)$ is $\OO(x^2)$ at infinity and
$\OO(x)$ at zero.

Order relations cannot in general be exponentiated:
$2x=\ord(x)$ as $x\to\infty$, whereas
$e^{2x}\ne\ord(e^x)$. Nor can they be differentiated without further
assumptions. For instance $x\sin x=\OO(x)$, but
$x\cos x+\sin x$ is not $\OO(1)$ at infinity.
Finally two functions of order $g$ may cancel, so their sum need not
be of order $g$.
\edit{Correction (p.~6)}{The source writes $x\sin x=\ord(x)$;
this fails because $\sin x$ has infinitely many zeros. The intended
warning about differentiation is valid with $\OO(x)$.}

\subsection{Asymptotic equivalence and scales}
\src{6--7}
\begin{definition}
We write $f\sim g$ at a specified limit if $g$ is eventually nonzero
and $f/g\to1$.
\end{definition}
Then $f-g=o(g)$ and $f=\ord(g)$. The notation $f\sim0$ has no
meaning in this ratio definition.
For example, as $x\to\infty$,
\[
 7x^{5/2}+18x^2(\ln x)^3+84\sim7x^{5/2}.
\]
Stirling's formula is
\[
 n!\sim\sqrt{2\pi n}\left(\frac ne\right)^n.
\]
\begin{definition}
A sequence $(\delta_n)$ is an asymptotic scale at a specified limit if
$\delta_{n+1}=o(\delta_n)$ for every $n$.
\end{definition}
Examples are $x^n$ as $x\to0$; $x^{-a_n}$ as $x\to\infty$ when
$a_{n+1}>a_n$; and $e^{Ax}x^\beta x^{\gamma n}$ as
$x\to\infty$ when $\gamma<0$. Another scale at $0^+$ is
\[
 1,\ x,\ x^2\ln x,\ x^2,\ x^3(\ln x)^2,\ x^3\ln x,\ x^3,
 \ x^4(\ln x)^3,\ldots.
\]

\subsection{Asymptotic series and finite expansions}
\src{7--8}
\begin{definition}
For an asymptotic scale $(\delta_n)$, we write
\[
 f(x)\sim\sum_{n=1}^\infty a_n\delta_n(x)
\]
if for every fixed $N\ge1$,
\begin{equation}\label{eq:asymdef}
 f(x)-\sum_{n=1}^N a_n\delta_n(x)=o(\delta_N(x)).
\end{equation}
A finite expansion through $N$ requires this condition for
$1\le m\le N$, with $m$ in place of $N$.
\end{definition}
The coefficients are unique for the specified scale:
\[
 a_n=\lim\frac{f(x)-\sum_{k=1}^{n-1}a_k\delta_k(x)}{\delta_n(x)}.
\]
The empty sum for $n=1$ is zero. An asymptotic series need not
converge, and it need not uniquely determine the function. As $x\to0^+$,
\[
 e^{-1/x}\sim\sum_{n=0}^\infty0\,x^n,
 \qquad e^x+e^{-1/x}\sim\sum_{n=0}^\infty\frac{x^n}{n!}.
\]
\edit{Correction (p.~7)}{The source calls a remainder
$\OO(\delta_{N+1})$ an alternative to \eqref{eq:asymdef}.
For an arbitrary finite expansion this is stronger, not equivalent.
For an infinite expansion satisfying the definition at \emph{every}
order, the next-order condition does imply it. The flat-function
example requires $x\to0^+$.}

A function may have expansions in different scales. For example,
\begin{align*}
 \tan x&=x+\frac{x^3}3+\frac{2x^5}{15}+\OO(x^7),\\
 &=\sin x+\frac12\sin^3x+\frac38\sin^5x+\OO(\sin^7x)
 \qquad(x\to0).
\end{align*}
Expansions may be added, subtracted, multiplied and (with a suitable
nonvanishing leading denominator) divided, retaining the appropriate
remainder orders and, if needed, enlarging the scale.

Substitution requires care. Let $f(x)=e^{x^2}$ and
$x(\eps)=\eps^{-1}+\eps$. Although $x(\eps)\sim\eps^{-1}$,
\[
 f(x(\eps))=e^{\eps^{-2}+2+\eps^2}\sim e^2e^{\eps^{-2}},
 \qquad f(\eps^{-1})=e^{\eps^{-2}}.
\]

\subsection{A fractional-power example}
\src{8}
Find two terms of
\[
 F(\eps)=\frac{[\ln(1+\eps)]^{1/4}}{\cosh(\eps^{1/3})},
 \qquad\eps\to0^+.
\]
Using
\begin{align*}
 \ln(1+\eps)&=\eps-\frac{\eps^2}2+\OO(\eps^3),\\
 (1+z)^\alpha&=1+\alpha z+\OO(z^2),\\
 \cosh z&=1+\frac{z^2}2+\OO(z^4),
\end{align*}
we obtain
\begin{align*}
 [\ln(1+\eps)]^{1/4}
 &=\eps^{1/4}\left(1-\frac\eps8+\OO(\eps^2)\right),\\
 \frac1{\cosh(\eps^{1/3})}
 &=1-\frac12\eps^{2/3}+\OO(\eps^{4/3}),\\
 F(\eps)&=\eps^{1/4}-\frac12\eps^{11/12}
 +\OO(\eps^{5/4}).
\end{align*}
\edit{Correction (p.~8)}{The generic binomial expansion has remainder
$\OO(z^2)$, not $\OO(z^3)$ after just two terms. The final remainder
is $\OO(\eps^{5/4})$; the handwritten exponent appears to be $5/24$.
The omitted third term is $-\eps^{5/4}/8$.}

\section{Perturbation of algebraic equations}
\src{9--10}
\subsection{Two regular roots}
Consider
\[
 (1-\eps)x^2+2x+\eps=0.
\]
At $\eps=0$ the roots are $0,-2$. For the root near zero,
\[
 x=-\frac\eps2-\frac12(1-\eps)x^2
 =-\frac\eps2+\OO(\eps^2).
\]
For the second root write $x=-2+\bar x$. Substitution gives
\[
 -2\bar x-3\eps+\OO(\bar x^2,\eps\bar x)=0,
 \qquad\bar x=-\frac32\eps+\OO(\eps^2).
\]
Hence
\begin{equation}
 x^{(1)}=-\frac\eps2+\OO(\eps^2),\qquad
 x^{(2)}=-2-\frac32\eps+\OO(\eps^2).
\end{equation}
\edit{Correction (p.~9)}{The source gives $-2\eps/3$ for the second
correction. Solving its own equation $-2\bar x-3\eps=0$ gives
$-3\eps/2$. Both roots admit ordinary power-series expansions locally.}

\subsection{A root lost in the limiting equation}
Now consider
\begin{equation}\label{eq:singquad}
 \eps x^2+x-1=0.
\end{equation}
Setting $\eps=0$ gives only $x=1$. Write $x=1+\bar x$,
then $\bar x\sim-\eps$. Putting $\bar x=-\eps+\widetilde x$ gives
\[
 \widetilde x=2\eps^2+\OO(\eps^3),\qquad
 x_+=1-\eps+2\eps^2+\OO(\eps^3).
\]
To find the other root, compare possible dominant balances.
Balancing $\eps x^2$ with $1$ suggests $x=\ord(\eps^{-1/2})$,
but then $x$ is larger than the supposedly dominant terms, so this
balance is inconsistent. Balancing $\eps x^2$ with $x$ instead gives
$x\sim-\eps^{-1}$, with both terms much larger than $1$.
Set $x=-\eps^{-1}+\bar x$. Then
\[
 -\bar x-1+\eps\bar x^2=0,\qquad
 \bar x=-1+\OO(\eps).
\]
Rescaling $x=X/\eps$ gives the regular equation
$X^2+X-\eps=0$, with $X\sim-1$ on this branch.
\edit{Correction and reconstruction (pp.~9--10)}{One line in the
source loses the minus sign in $x=-1/\eps+\bar x$.
Using the sum of the two roots in \eqref{eq:singquad} completes the
second expansion:
$x_-=-\eps^{-1}-1+\eps-2\eps^2+\OO(\eps^3)$.}

\section{Transcendental equations}
\subsection{An exponentially small correction}
\src{10}
Solve $x+\tanh x=\lambda$ as $\lambda\to\infty$.
First $x=\lambda+\bar x$ with $\bar x\sim-1$. Writing
$x=\lambda-1+\widetilde x$ and using
\[
 \tanh x=\frac{1-e^{-2x}}{1+e^{-2x}}
 =1-2e^{-2x}+\OO(e^{-4x})
\]
gives $\widetilde x\sim2e^{-2(\lambda-1)}$. Thus
\[
 x=\lambda-1+2e^{-2(\lambda-1)}
 +\OO(e^{-4(\lambda-1)}).
\]
\edit{Reconstruction (p.~10)}{The explicit remainder follows by
substitution and $\widetilde x=\OO(e^{-2(\lambda-1)})$.
Approximating the argument of $\tanh$ is justified at this
exponentially small scale, not merely by $\bar x\sim-1$.}

\subsection{\texorpdfstring{Two positive roots: $x^2-\ln x+1=\lambda$}{Two positive roots: logarithmic equation}}
\src{11--12}
The source's sketch shows a curve tending to infinity at both ends
of $x>0$, with a minimum between them. Indeed
$f'(x)=2x-1/x$ has its only zero at $x=1/\sqrt2$;
there are exactly two positive roots for sufficiently large $\lambda$.

For the large root, write $x=\sqrt\lambda+\bar x$,
$\bar x=o(\sqrt\lambda)$. Then
\[
 2\sqrt\lambda\,\bar x+\bar x^2
 -\frac12\ln\lambda
 -\ln\left(1+\frac{\bar x}{\sqrt\lambda}\right)+1=0.
\]
The leading correction is $\bar x\sim\ln\lambda/(4\sqrt\lambda)$.
Taking the constant residual into account gives
\begin{equation}
 x_{\mathrm{large}}=\sqrt\lambda
 +\frac{\ln\lambda-2}{4\sqrt\lambda}
 +\OO\left(\frac{(\ln\lambda)^2}{\lambda^{3/2}}\right).
\end{equation}
One may instead display the second term as two separate terms
$\ln\lambda/(4\sqrt\lambda)-1/(2\sqrt\lambda)$.

For the small root, $-\ln x\sim\lambda$ does not by itself imply
$x\sim e^{-\lambda}$. Put $x=e^s$ to obtain
\[
 e^{2s}-s+1=\lambda.
\]
Let $s=-\lambda+\bar s$, so $\bar s\sim1$, and then
$\bar s=1+\widetilde s$. The exact reduced equation is
\[
 \widetilde s=e^{2-2\lambda}e^{2\widetilde s}.
\]
Therefore, with $q=e^{2-2\lambda}$,
\[
 \widetilde s=q+\OO(q^2),\qquad
 x_{\mathrm{small}}=e^{1-\lambda}[1+q+\OO(q^2)].
\]
\edit{Reconstruction (pp.~11--12)}{Remainders are supplied for both
branches; the substitution into the logarithmic variable avoids
unjustified exponentiation of an asymptotic equivalence.}

\subsection{\texorpdfstring{Two roots: $xe^{1/x}=e^\lambda$}{Two roots: exponential equation}}
\src{12}
On $x>0$, the function decreases to its minimum $e$ at $x=1$ and
then increases; hence there are two roots for $\lambda>1$.
For the large root, put $E=e^\lambda$ and $x=E+x_1$.
Expansion of $xe^{1/x}=x+1+1/(2x)+\OO(x^{-2})$ gives
$x_1\sim-1$. Taking $x=E-1+x_2$ then gives
\[
 x_2\sim-\frac1{2E},\qquad
 x_{\mathrm{large}}=e^\lambda-1-\frac12e^{-\lambda}
 +\OO(e^{-2\lambda}).
\]
For the small root, take logarithms:
\[
 \ln x+\frac1x=\lambda,\qquad x\sim\frac1\lambda.
\]
Write $x=\lambda^{-1}+x_1$, $x_1=o(\lambda^{-1})$. Then
\[
 \ln x=-\ln\lambda+o(1),\quad
 \frac1x=\lambda-\lambda^2x_1+\OO(\lambda^3x_1^2),
\]
so
\[
 x_1\sim-\frac{\ln\lambda}{\lambda^2},\qquad
 x_{\mathrm{small}}=\frac1\lambda-
 \frac{\ln\lambda}{\lambda^2}
 +\OO\left(\frac{(\ln\lambda)^2}{\lambda^3}\right).
\]
\edit{Correction (p.~12)}{The source's final correction for the small
root is positive, but $\ln(1/\lambda)=-\ln\lambda$ makes it negative.
One intermediate large-root line likewise has the wrong sign for
$x_1$; the source's final large-root answer has the correct minus sign.}

\subsection{\texorpdfstring{Large roots of $x\tan x=1$}{Large roots of x tan x = 1}}
\src{13}
Since $\tan x=1/x$ is small for a large positive root, write
$x=n\pi+h_n$, $0<h_n<\pi/2$. Tangent is $\pi$-periodic, so
\[
 (n\pi+h_n)\tan h_n=1.
\]
The dominant balance is $n\pi h_n\sim1$, giving
\[
 x_n=n\pi+\frac1{n\pi}+\OO(n^{-3})\qquad(n\to\infty).
\]
\edit{Correction and reconstruction (p.~13)}{The source assigns a
minus sign to $\tan(n\pi+h)$ when $n$ is odd. In fact
$\tan(n\pi+h)=\tan h$ for every integer $n$.
The remainder is obtained from $\tan h=h+h^3/3+\OO(h^5)$;
more precisely the next correction is $-4/[3(n\pi)^3]$.}

\subsection{An exponentially narrow hyperbolic term}
\src{13}
Consider
\[
 x+1+\frac1{\cosh(x/\eps)}=0,\qquad\eps\to0^+.
\]
The source sketch compares the narrow positive peak of
$1/\cosh(x/\eps)$ with the line $-x-1$.
Any root satisfies $x<-1$; on this interval the sum on the left is
strictly increasing, so there is exactly one root.
Set $x=-1+h$, where $h<0$ is small. Then
\[
 \frac1{\cosh((-1+h)/\eps)}
 \sim2e^{-1/\eps}e^{h/\eps},\qquad
 h\sim-2e^{-1/\eps}.
\]
Thus
\[
 x=-1-2e^{-1/\eps}
 +\OO\left(\frac{e^{-2/\eps}}\eps\right).
\]
\edit{Reconstruction (p.~13)}{The small correction satisfies
$h/\eps\to0$, which justifies expanding $e^{h/\eps}$.
Uniqueness and the remainder are added.}

\section{Local analysis of ordinary differential equations}
\src{14--16}
Consider
\begin{equation}\label{eq:ode}
 y''+p(x)y'+q(x)y=0
\end{equation}
near $x=x_0$, using $x-x_0$ as the small variable.
\subsection{Ordinary points and power series}
An ordinary point is one at which $p$ and $q$ are analytic.
At such a point solutions have convergent power-series representations
locally. For $x_0=0$, seek $y=\sum_{n=0}^\infty a_nx^n$.
\edit{Clarification (p.~14)}{The source defines a regular point using
finite limits of $p,q$. Continuity suffices for a local existence theory,
but analyticity is required for the power-series conclusion used here.}

For $y''+y=0$, equating coefficients gives
\[
 (n+2)(n+1)a_{n+2}+a_n=0,\quad
 a_{2k}=\frac{(-1)^ka_0}{(2k)!},\quad
 a_{2k+1}=\frac{(-1)^ka_1}{(2k+1)!}.
\]
Consequently
\[
 y=a_0\sum_{k=0}^\infty\frac{(-1)^kx^{2k}}{(2k)!}
 +a_1\sum_{k=0}^\infty\frac{(-1)^kx^{2k+1}}{(2k+1)!}
 =a_0\cos x+a_1\sin x.
\]

\subsection{The Airy equation}
\src{14--15}
For $y''-xy=0$, power-series substitution gives
\[
 2a_2=0,\qquad(n+3)(n+2)a_{n+3}=a_n\quad(n\ge0).
\]
Thus $a_2=a_5=a_8=\cdots=0$, and
\begin{align*}
 a_{3n}&=\frac{a_0}{(3n)(3n-1)(3n-3)(3n-4)\cdots3\cdot2}
 =\frac{a_0\Gamma(2/3)}{3^{2n}n!\Gamma(n+2/3)},\\
 a_{3n+1}&=\frac{a_1\Gamma(4/3)}{3^{2n}n!\Gamma(n+4/3)}.
\end{align*}
Define the normalized solutions
\begin{align}
 y_1(x)&=\Gamma(2/3)\sum_{n=0}^\infty
 \frac{x^{3n}}{3^{2n}n!\Gamma(n+2/3)},\\
 y_2(x)&=\Gamma(4/3)\sum_{n=0}^\infty
 \frac{x^{3n+1}}{3^{2n}n!\Gamma(n+4/3)}.
\end{align}
Then $y=a_0y_1+a_1y_2$, with
$y_1(0)=1,y_1'(0)=0,y_2(0)=0,y_2'(0)=1$.
\edit{Ambiguous source normalization (p.~15)}{The handwritten underbraces
under the two series leave unclear whether the Gamma prefactors are
part of $y_1,y_2$. The source then writes
$\operatorname{Ai}(x)=3^{-2/3}y_1(x)-3^{-4/3}y_2(x)$ and
$\operatorname{Bi}(x)=3^{-1/6}y_1(x)+3^{-5/6}y_2(x)$.
These are transcribed here for traceability, but should \emph{not}
be used as standard Airy normalizations with our explicitly normalized
$y_1,y_2$. The power-series solutions themselves are unambiguous.}
\edit{Verified reconstruction}{Using the initial values in
\href{https://dlmf.nist.gov/9.2.ii}{NIST DLMF, Section 9.2(ii)},
the standard normalizations for our basis are}
\begin{align*}
 \operatorname{Ai}(x)&=\frac{y_1(x)}{3^{2/3}\Gamma(2/3)}
 -\frac{y_2(x)}{3^{1/3}\Gamma(1/3)},\\
 \operatorname{Bi}(x)&=\frac{y_1(x)}{3^{1/6}\Gamma(2/3)}
 +\frac{3^{1/6}y_2(x)}{\Gamma(1/3)}.
\end{align*}

\subsection{Regular singular points and Euler equations}
\src{15--16}
At a regular singular point, $(x-x_0)p(x)$ and
$(x-x_0)^2q(x)$ extend analytically to $x_0$.
For $4x^2y''+y=0$ at zero, an ordinary power series would require
\[
 [4n(n-1)+1]a_n=(2n-1)^2a_n=0.
\]
For integer $n\ge0$ this forces every coefficient to vanish, so
an ordinary Taylor series cannot represent a nonzero solution.

For the general Euler equation
\[
 x^2y''+\alpha xy'+\beta y=0,\qquad x>0,
\]
try $y=x^r$. The indicial equation is
\[
 r(r-1)+\alpha r+\beta=0.
\]
Distinct real roots give $x^{r_1},x^{r_2}$. For a repeated root $r$,
\[
 \lim_{r_2\to r}\frac{x^{r_2}-x^r}{r_2-r}=x^r\ln x,
\]
and the fundamental set is $x^r,x^r\ln x$.
In particular $4x^2y''+y=0$ has solutions
$x^{1/2}$ and $x^{1/2}\ln x$.
\edit{Clarification (pp.~15--16)}{Finite weighted limits alone are
weaker than the usual analytic hypotheses for Frobenius theory.
The logarithmic example is valid on $x>0$.}

\subsection{Frobenius series}
For a regular singular point at zero, seek
\[
 y=x^r\sum_{n=0}^\infty a_nx^n,\qquad a_0\ne0.
\]
Consider $2x^2y''-xy'+(1+x)y=0$. The indicial equation is
\[
 2r(r-1)-r+1=(2r-1)(r-1)=0,
 \qquad r=1,\tfrac12.
\]
For $n\ge1$,
\[
 [2(n+r)-1][n+r-1]a_n+a_{n-1}=0.
\]
For $r=1$,
\[
 a_n=-\frac{a_{n-1}}{n(2n+1)}
 =\frac{(-1)^na_0\,2^n}{(2n+1)!},\qquad
 y_{(1)}=x\sum_{n=0}^\infty\frac{(-1)^n2^nx^n}{(2n+1)!}.
\]
\edit{Reconstruction (p.~16)}{The notes say to find the other branch
but do not work it out. For $r=1/2$,
$a_n=-a_{n-1}/[n(2n-1)]$, giving
$y_{(2)}=x^{1/2}\sum_{n=0}^\infty(-1)^n2^nx^n/(2n)!$.
On $x>0$ these are respectively
$\sqrt{x/2}\sin\sqrt{2x}$ and $\sqrt{x}\cos\sqrt{2x}$;
arbitrary constants multiply the two branches.}

\section{An irregular singular point: Bessel behaviour at infinity}
\src{16--18}
Consider the order-zero Bessel equation
\begin{equation}\label{eq:bessel}
 y''+\frac1xy'+y=0.
\end{equation}
Zero is regular singular, while infinity is irregular singular
(after changing variable to $t=1/x$). Dropping $y'/x$ suggests
oscillations $e^{\pm ix}$, but misses their algebraic amplitude.
\edit{Clarification (pp.~16--17)}{The source's warning that
$xe^{\pm ix}$ also has the same leading derivative balance is
understood as a formal balance, not as an exact solution of
$y''+y=0$ or of \eqref{eq:bessel}.}
For an individual complex oscillatory branch, put $y=e^{S(x)}$.
Then
\[
 S''+(S')^2+\frac{S'}x+1=0.
\]
Let $S'=\sigma i+T$, $\sigma\in\{1,-1\}$. It follows that
\[
 T'+2\sigma iT+T^2+\frac{\sigma i}x+\frac Tx=0.
\]
Assuming $T'=o(T)$ and $T=o(1)$ gives the consistent balance
$T\sim-1/(2x)$. Iterating gives
\[
 T=-\frac1{2x}+\OO(x^{-2}),\qquad
 S=\sigma ix-\frac12\ln x+C_\sigma+\OO(x^{-1}),
\]
and therefore the two oscillatory modes have the form
\[
 y_\sigma=C_\sigma x^{-1/2}e^{\sigma ix}
 [1+\OO(x^{-1})].
\]
\edit{Clarification (p.~17)}{This is a formal asymptotic derivation;
the assumptions on $T$ are self-consistent, but are not by themselves
a full existence proof. A general real solution is a combination of
the two modes and has zeros, so a global logarithm $S$ need not exist.}

\subsection{Near zero and the standard Bessel solutions}
Near zero the regular solution and an independent logarithmic solution
can be normalized as
\begin{align*}
 J_0(x)&=1-\frac{x^2}4+\OO(x^4),\\
 L_0(x)&=\ln x-\frac{x^2}4\ln x+\frac{x^2}4
 +\OO\bigl(x^4|\ln x|\bigr)\qquad(x\to0^+).
\end{align*}
The standard second-kind normalization is
\[
 Y_0(x)=\frac2\pi\bigl[L_0(x)+(\gamma-\ln2)J_0(x)\bigr].
\]
The large-$x$ leading oscillations are conventionally written
\begin{align*}
 J_0(x)&=\sqrt{\frac2{\pi x}}
 \left[\cos\left(x-\frac\pi4\right)+\OO(x^{-1})\right],\\
 Y_0(x)&=\sqrt{\frac2{\pi x}}
 \left[\sin\left(x-\frac\pi4\right)+\OO(x^{-1})\right].
\end{align*}
\edit{Correction (p.~18)}{The source gives cosine for both $J_0$
and $Y_0$; the second phase is sine. Additive error notation is used
because ratio equivalence is unsuitable at zeros of the oscillation.
The handwritten instruction assigning constants to the logarithmic
solution is ambiguous: $Y_0$ requires adding a multiple of $J_0$ as
well as scaling $L_0$. The phase and normalization constants are
stated here, as in the source, without deriving their global matching.
The normalizations were checked against
\href{https://dlmf.nist.gov/10.8}{NIST DLMF, Section 10.8}, and the
phases against \href{https://dlmf.nist.gov/10.17}{Section 10.17}.}

\section{Solvability, resonance and eigenvalue perturbation}
\subsection{Inner products and adjoints}
\src{18--19}
For $x,y\in\mathbb C^n$, use the convention
\[
 \langle x,y\rangle=x^T\overline y,
\]
linear in the first argument and conjugate-linear in the second.
Then $\langle x,x\rangle\ge0$, with equality precisely when $x=0$.
The adjoint satisfies
\[
 \langle Ax,y\rangle=\langle x,A^*y\rangle.
\]
It is unique: if $W_1,W_2$ both satisfy this identity, then
$\langle x,(W_1-W_2)y\rangle=0$ for all $x,y$. Taking
$x=(W_1-W_2)y$ proves $W_1=W_2$.
Direct calculation gives $A^*=\overline A^{\,T}$.
If $A=A^*$, it is self-adjoint and all its eigenvalues are real:
if $Av=\lambda v$ and $v\ne0$, then
\[
 \lambda\langle v,v\rangle
 =\langle Av,v\rangle
 =\langle v,Av\rangle
 =\overline\lambda\langle v,v\rangle.
\]

\subsection{The Fredholm alternative in finite dimensions}
An invertible (regular) matrix has only the trivial null vector and
$Ax=b$ has a unique solution for every $b$. If $A$ is singular,
the system is either inconsistent or has infinitely many solutions.
\edit{Correction (p.~19)}{The source's definition of regularity must
read ``only the trivial solution'' of $Ax=0$; the zero vector is
always a solution, even for a singular matrix.}

If $Ax=b$ and $v\in\Null(A^*)$, then
\[
 \langle b,v\rangle=\langle Ax,v\rangle
 =\langle x,A^*v\rangle=0.
\]
In fact this necessary condition is sufficient:
\begin{equation}\label{eq:fredholm}
 Ax=b\text{ is solvable}\quad\Longleftrightarrow\quad
 \langle b,v\rangle=0\quad\text{for all }v\in\Null(A^*).
\end{equation}
If $v^{(1)},\ldots,v^{(m)}$ form a basis of $\Null(A^*)$, it suffices
to test these $m$ vectors. Since $\rank A=\rank A^*$ and
rank--nullity holds, $\dim\Null A=\dim\Null A^*$ for square matrices.
\edit{Reconstruction (p.~19)}{For sufficiency, the identity
$(\operatorname{Ran} A)^\perp=\Null(A^*)$ implies
$\operatorname{Ran} A=\Null(A^*)^\perp$ in finite dimensions.}

\subsection{A regular matrix perturbation}
\src{19}
Consider
\[
 (A+\eps B)x=b.
\]
If $A$ is invertible, substitute $x\sim x_0+\eps x_1+\cdots$:
\[
 Ax_0=b,\quad Ax_1+Bx_0=0,\qquad
 x_0=A^{-1}b,\quad x_1=-A^{-1}BA^{-1}b.
\]
This expansion generally fails when $A$ is singular.

\subsection{Simple resonance}
\src{19--20}
Assume $\Null A=\spanof\{u\}$ and $\Null A^*=\spanof\{v\}$.
Try a Laurent expansion
\[
 x\sim\eps^{-1}x_{-1}+x_0+\eps x_1+\cdots.
\]
Successive orders give
\begin{align*}
 \eps^{-1}:&\quad Ax_{-1}=0,\quad x_{-1}=au,\\
 1:&\quad Ax_0=b-aBu.
\end{align*}
By \eqref{eq:fredholm},
\[
 \langle b,v\rangle-a\langle Bu,v\rangle=0.
\]
If $\langle Bu,v\rangle\ne0$, then
\begin{equation}
 a=\frac{\langle b,v\rangle}{\langle Bu,v\rangle},\qquad
 x_{-1}=\frac{\langle b,v\rangle}{\langle Bu,v\rangle}u.
\end{equation}
Let $p$ be any particular solution of $Ap=b-aBu$. Then
$x_0=p+a'u$. The next equation $Ax_1=-Bx_0$ fixes
\[
 a'=-\frac{\langle Bp,v\rangle}{\langle Bu,v\rangle}.
\]
If $\langle b,v\rangle=0$, the leading pole coefficient vanishes
and the unperturbed equation is solvable.

\subsection{Super-resonance}
\src{20}
If $\langle Bu,v\rangle=0$, the preceding quotient is unavailable.
Try instead
\[
 x\sim\eps^{-2}x_{-2}+\eps^{-1}x_{-1}+x_0+\cdots.
\]
At the first two orders,
\[
 x_{-2}=au,\qquad Ax_{-1}=-aBu.
\]
Since $Bu$ is orthogonal to $\Null(A^*)$, choose a particular
solution $p$ of $Ap=-Bu$. Then $x_{-1}=ap+a'u$. At order one,
\[
 Ax_0=b-B(ap+a'u).
\]
The solvability condition is
\[
 \langle b,v\rangle-a\langle Bp,v\rangle=0,
 \qquad a=\frac{\langle b,v\rangle}{\langle Bp,v\rangle}
 \quad\text{if }\langle Bp,v\rangle\ne0.
\]
\edit{Clarification (p.~20)}{The definition $Ap=-Bu$, implicit in the
handwritten change of notation, is made explicit. A second-order pole
is conditional on the new nonzero denominator; if it also vanishes,
further analysis is required.}

\subsection{Eigenvalue perturbation}
\src{20--21}
Seek nonzero solutions of
\[
 (A+\eps B)x=\lambda x.
\]
Let $a$ be an algebraically simple eigenvalue of $A$, and write
\[
 A_a=A-aI,\quad\Null A_a=\spanof\{u\},\quad
 \Null A_a^*=\spanof\{v\}.
\]
Expand
\[
 \lambda=a+\eps\lambda_1+\eps^2\lambda_2+\cdots,
 \qquad x=x_0+\eps x_1+\cdots.
\]
At leading order, $Ax_0=ax_0$, so $x_0=ku$, $k\ne0$.
At the next order,
\[
 A_ax_1=-Bx_0+\lambda_1x_0=-kBu+k\lambda_1u.
\]
The Fredholm condition yields
\begin{equation}
 \boxed{\lambda_1=\frac{\langle Bu,v\rangle}{\langle u,v\rangle}}.
\end{equation}
It is independent of the normalization $k$. For self-adjoint $A$,
$a$ is real and one may take $v=u$, so
\[
 \lambda_1=\frac{\langle Bu,u\rangle}{\langle u,u\rangle}.
\]
\edit{Clarification (p.~21)}{A one-dimensional eigenspace alone does
not guarantee a simple eigenvalue. For an algebraically simple
eigenvalue, $\langle u,v\rangle\ne0$. If this pairing vanishes,
the eigenvalue is defective and the assumed ordinary power series
may fail. The source notes this possible zero denominator but does
not resolve it. Eigenvector normalization is still needed to make
$x_1$ unique.}

\appendix
\section{Editorial register and remaining uncertainties}
All 21 source pages were read visually. The following list records
substantive departures from the handwriting; minor spelling,
punctuation, dummy-index changes, and routine alignment are not listed.
\begingroup\small
\setlength{\parskip}{2pt}
\begin{enumerate}
\setlength{\itemsep}{2pt}
 \item Page 1: corrected the logarithmic inequality's sign and
 qualified the strict harmonic-number bounds at $n=1$.
 \item Page 2: corrected first-order exponential Taylor signs and
 supplied a justification for the limiting integral for $\gamma$.
 \item Pages 3--4: supplied the binomial series convergence domain
 and the exclusion of Gamma poles from finite-valued continuation.
 \item Page 5: preserved the course's $\ord$ convention; clarified
 one-sided logarithmic limits and the unusual $x^2-1$ example.
 \item Page 6: replaced the false $x\sin x=\ord(x)$ claim by
 $x\sin x=\OO(x)$ in the differentiation warning.
 \item Page 7: distinguished finite-expansion remainder conditions;
 specified $0^+$ for the flat-function example.
 \item Page 8: corrected the generic binomial remainder and the
 final fractional-power remainder exponent.
 \item Page 9: changed the root correction $-2\eps/3$ to
 $-3\eps/2$ and restored a missing minus sign in the large-root ansatz.
 \item Pages 10--13: supplied remainders and consistency arguments
 in the transcendental examples. These are deductions, not literal
 transcriptions of additional handwritten terms.
 \item Page 12: corrected the small-root logarithmic sign and a
 conflicting intermediate sign in the large-root calculation.
 \item Page 13: corrected tangent periodicity for odd $n$; completed
 the large-root error estimate and the hyperbolic-root uniqueness check.
 \item Pages 14--16: stated the analytic hypotheses for power-series
 and Frobenius theory; completed the second Frobenius branch.
 \item Page 15: the Airy $y_1,y_2$ underbrace normalization remains
 ambiguous in the source. Its named-function combinations are retained
 inside a warning; verified standard normalizations are supplied
 separately for the explicitly defined basis.
 \item Pages 16--17: distinguished formal dominant balance from
 exact solutions and from a proof of asymptotic existence.
 \item Page 18: corrected the $Y_0$ phase, expressed its near-zero
 normalization as a linear combination, and used additive errors
 for oscillatory asymptotic formulas.
 \item Page 19: corrected the definition of matrix regularity and
 supplied the range--nullspace proof of sufficiency.
 \item Page 20: made the super-resonance particular solution and
 nondegeneracy requirement explicit.
 \item Page 21: distinguished algebraic simplicity from a
 one-dimensional eigenspace and stated the normalization limitation.
\end{enumerate}
The three qualitative source sketches (two-root logarithmic equation,
two-root exponential equation, and narrow hyperbolic peak) are
described in the corresponding text rather than reproduced as images.
No mathematical content is inferred from their drawing scale.
\endgroup
\end{document}
